En la vida quotidiana estem sotmesos a moviments vibratoris. Per exemple, en caminar, córrer, viatjar amb algun mitjà de locomoció o estar a prop d'alguna màquina. A l'hora de dissenyar vehicles i màquines, cal fer un estudi d'aquests moviments per tal d'aconseguir que siguin confortables i segurs, ja que els efectes de les vibracions poden anar des de simples molèsties fins al dolor o la mort.

Aquests estudis solen utilitzar l'acceleració màxima del moviment vibratori com a variable, per a relacionar-la amb les molèsties que percebem.

Se sap que som molt sensibles a un moviment vibratori de 6,0~Hz i que, amb aquesta freqüència, a partir d'una acceleració màxima de $6,0\ \mathrm{m\,s^{-2}}$, les molèsties són tan fortes que ens poden arribar a alarmar.

\begin{apartat}{1}
Calculeu l'amplitud d'oscil·lació que correspon a un moviment vibratori harmònic de 6,0~Hz i una acceleració màxima de $6,0\ \mathrm{m\,s^{-2}}$.
\end{apartat}

\begin{apartat}{1}
Calculeu el valor de la constant elàstica d'una molla per tal que una massa de 85~kg que hi estigui enganxada oscil·li amb una freqüència de 6,0~Hz.
\end{apartat}
